\subsection{IF-11: an inhabited two-step algebraic theta-jet specialization}
The numerical examples in IF-09 are not the only source of coefficient
data. A separate supplied construction gives a genuine theta family
with $m=2$ and all the algebraicity needed by its finite kernels. This
does not assert the additional global size or nonzero-period hypotheses.

All paths in this subsection beginning with \texttt{Development/} are
paths in Alp\"oge's Lean archive. Short paths beginning with \texttt{Proofs/}
or \texttt{Definitions/} below are relative to
\texttt{Development/AbelianVarietyTheory/}. The theorem
\texttt{GlobalTheory.KerT.JetsInhab.exists\_colIn\_jets}, in
\texttt{Development/AbelianVarietyTheory/Proofs/GlobalInterpolation/}
\texttt{ExistsHypJetsAll.lean}, lines 73--115, supplies
\[
 \exists B,X,I,\quad
 \texttt{HasLetters}(B)\ \wedge\ I.\texttt{Ok}\ \wedge\ I.\texttt{JetsIn}.
                                                               \tag{IF-11a}
\]
We use this as a cited theorem. Its proof builds its base from
\texttt{ExistsNatBoundPre.lean}, \texttt{exists\_pre\_two\_words}
(lines 29--132), and then fixes the distinguished word in
\texttt{exists\_pre\_sized} (lines 19--71 of
\texttt{ExistsHypJetsAll.lean}). Its rank datum is
\texttt{Formalism/SharpRestriction/InertJacobiLevels/HostFam.lean},
\texttt{Inhab.Rt 0} (lines 14--29). Thus
\[
 k=8,\quad m=2,\quad t=15,\quad f=8,\qquad
 (N_0,N_1,N_2)=(8,24,40).
\]

\paragraph{Statement IF-11.}
The witness constructed in (IF-11a) has the following coefficient
specialization, constant over its eight embedding factors:
\[
\begin{split}
 u&=(49,45,37,29,21,13,5,-1)/2,\\
 v&=(24,20,16,12,8,4,0),\\
 b&=(1,37),\qquad e=(19,55),\\
 \lambda_0&=(21,20,17,14,11,8,5,3).
\end{split}                                                     \tag{IF-11b}
\]
Its actual column determinant exponents and polynomial parts are
\[
 (s_0,s_1,s_2)=(3,-15,-43),
\]
\[
\begin{split}
 \alpha_0&=(18,17,14,11,8,5,2,0),\\
 \alpha_1&=(28,27,24,21,18,15,12,10,0^{[16]}),\\
 \alpha_2&=(48,47,44,41,38,35,32,30,20^{[16]},0^{[16]}).
\end{split}                                                     \tag{IF-11c}
\]
The two finite kernels have shapes
\[
 \nu_1=(28,27,24,21,18,15,12,10),\qquad
 \nu_2=(48,47,44,41,38,35,32,30,20^{[16]}),          \tag{IF-11d}
\]
with determinant insertions $(10,20)$, of degrees $80,480$ per
embedding. All required normalized moments lie in the finite number
field supplied by (IF-11a). Hence IF-08 and IF-10 apply to this actual
theta-jet family, not only to the exponential test series.

The distinct nonzero simple-root gaps in these three coefficient
lattices are $1,2,3,10,20$. The gap $1$ has no second descendant; for
the others the exact second-descendant invariants are
\[
\begin{array}{c|c|c|c}
 d&F^2\text{ on Weyl}&F^2\text{ on normalized dual}
       &F^{(2)}\text{ on normalized dual}\\\hline
 2&2&2&1\\
 3&2&6&3\\
 10&2&90&45=3^2\cdot5\\
 20&2&380&190=2\cdot5\cdot19
\end{array}.                                                    \tag{IF-11e}
\]
On every nonzero displayed Weyl second-descendant line,
$F^{(2)}$ has invariant one. These are coefficient-root indices,
not the values or Smith invariants of the full global theta composition.

\begin{proof}
The parameter construction
\texttt{Proofs/IwasawaDecomposition/ExistsHarishChandraWeights.lean},
\texttt{exists\_preCh\_of\_prescribed} (lines 12--33), gives
\[
 v_l=(k-2-l)4\Delta,\qquad
 b_j=b_1+2hj,\qquad e_j=b_j+h,
\]
and endpoint values $u_0=u_F+1/2$, $u_{k-1}=u_L+1/2$, with
$u_{l+1}=v_l-1/2-\Delta$ for the distinguished all-true word.
The invocation in \texttt{exists\_pre\_two\_words} has
\[
 \Delta=1,\quad h=t+3=18,\quad u_F=24,\quad u_L=-1,
 \quad b_1=1.
\]
Its distinguished word is $0\in\operatorname{Fin}2$, on which
\texttt{decide}(u=0) is true. These substitutions prove (IF-11b).
The passage to the one-word base keeps this weight family. The field
enlargement retains the preliminary rank and weight data, and the
subsequent change of ambient theta data retains the resulting base $B$. Thus the actual witness in (IF-11a) retains these
parameters; no selection of an unrelated existential witness is used.

Subtracting the rank-eight half-sum of positive roots gives the
displayed $\lambda_0$. From \texttt{Params.wG} one obtains
$w_G=(7,27)$, so the unshifted weights are
\[
 \lambda^N_1=(\lambda_0,-7^{[16]}),\qquad
 \lambda^N_2=(\lambda_0,-7^{[16]},-27^{[16]}).
\]
Subtract $8j$ at stage $j$, as required by
\texttt{Wts.lamTrue}, and then subtract the last entry to obtain
(IF-11c). Here
\[
 \mathrm{sSharp}=(7,35),\qquad
 \mathrm{sDiv}=(15,43),\qquad
 n_j=s_j+\mathrm{sSharp}_j=(10,20).
\]
Adding these $n_j$ to each source polynomial part gives (IF-11d),
which also agrees by direct substitution with \texttt{nuRow}.
The determinant degrees are $8\cdot10=80$ and $24\cdot20=480$.
Successive differences in (IF-11c) give the stated root gaps.
The invariant table is now exactly IF-04, with the specified Weyl
or normalized-dual lattice; no integral identification of those two
lattices is assumed.

It remains to check that the cited algebraicity is exactly the one
used by the kernel. In
\texttt{Definitions/AutomorphicMembership/JetsInIffStepField.lean},
\texttt{ColIn.jetNum} (lines 99--103) is
\[
 J(\beta)=\chi_F^{-1}
       \left(\prod_{\tau,u,i}c_{D,\tau}^{\beta_{\tau,u,i}}\right)
       (\partial^\beta\Theta)(0),                 \tag{IF-11f}
\]
where $\Theta$ is the actual \texttt{thInd} family. This is precisely
\texttt{ThDistC.jetC}, with ordinary derivatives and no division by
$\beta!$. Its \texttt{JetsIn} definition (lines 109--113) asserts
membership in the coefficient field for multidegree
$\sum_r\nu_{j,r}$ at each $\tau$, and
\texttt{JetsIn.exists\_eq,isAlgebraic} (lines 142--150) identifies
these numbers with their field-valued versions.

Every coefficient $A_\beta$ used by the polynomial kernel of shape
$\nu_j$ has exactly this multidegree: each exterior column contributes
its height, so the total is $|\nu_j|$ at that embedding. Consequently
the degree-restricted algebraicity hypothesis proves algebraicity of
every moment that actually occurs in the finite operator. An
all-degree hypothesis is unnecessary. The degrees here are
$155,635$ per embedding, or $1240,5080$ over the eight embeddings.

The field existence is not an assertion about the default-valued
function \texttt{desc}. Its supplied proof is
\texttt{Proofs/IwasawaDecomposition/DegreeEqTotDeg.lean},
\texttt{finAlg\_step} (lines 193--226), followed by
\texttt{exists\_jetsIn} (lines 283--313): a finite algebraic generating
set of the required theta jets is placed in the coefficient field,
using the supplied normalized-vacuum algebraicity theorem
\texttt{GlobalTheory.T.vacuum\_entire\_and\_normalizedJets\_algebraic}
in \texttt{VacuumEntireAndNormalizedJetsAlgebraicStatement.lean},
and theta equivariance and lattice reduction.
\texttt{ExistsColInJetsIn.exists\_bound\_jetsIn} (lines 34--40)
specifies the finite algebraic generating set to adjoin; the
\texttt{BaseExt.exists\_base} call in \texttt{exists\_colIn\_jets}
(lines 85--96) performs that adjunction while retaining the preliminary
rank and weight data.
Finally \texttt{Definitions/GlobalTwisting/JetsAllOfColIn.lean},
\texttt{jetsAll\_of\_colIn} (lines 108--116), identifies these
degree-restricted jets with the actual assembled kernel's
\texttt{JetsAll}. This proves exact compatibility with IF-08 and
discharges IF-10's coefficient-field input for this inhabited family.
The cited theorems are used as supplied results; we have not run a
new Lean kernel check of their proofs.
\end{proof}

\paragraph{A concrete derivative test and a full-module invariant.}
Write $\Delta_h(Z)$ for the leading $h$-row, $h$-column minor.
At one embedding factor the two highest flag polynomials for (IF-11d)
are the explicit polynomials
\[
 P_1=\Delta_1\Bigl(\prod_{h=2}^6\Delta_h^3\Bigr)
            \Delta_7^2\Delta_8^{10},\qquad
 P_2=P_1\Delta_{24}^{20}.                         \tag{IF-11g}
\]
For the full coefficient tensor put
$P_j^{\rm tot}=\prod_\tau P_j(Z_\tau)$. In standard coefficient
frames, before the scalar factors in (IF-08k), the normalized
highest-to-highest block coefficient is the moment functional
$J(P_j^{\rm tot})=\sum_\beta[Z^\beta]P_j^{\rm tot}\,J(\beta)$.
For other frames one uses the complete matrix conjugation (IF-08k);
one does not replace it by a scalar change of this one entry.
To prove the standard-frame formula,
apply the target highest functional to $A_{\nu_j}(Z)t_{\nu_j}$.
Cauchy--Binet selects the leading minors at each embedding, which gives
$P_j^{\rm tot}(Z)$; then contract its coefficients with $J$. Thus
(IF-11g) specifies an actual differential-operator test of degrees
$155,635$ per embedding, or $1240,5080$ in total. For the
coordinate exponential moments it equals one; no such value is
asserted for the actual analytic theta family.

The source coefficient representation has the wrong rank to be a
whole symplectic exterior algebra. The Weyl dimension
formula cited in IF-05 gives
\[
 \operatorname{rank}T_{\alpha_0}(8)
  =2^{28}\cdot3\cdot5^2\cdot7\cdot11^2\cdot19\cdot23.
                                                               \tag{IF-11h}
\]
This follows by multiplying
$\prod_{i<j}(\alpha_{0,i}-\alpha_{0,j}+j-i)/(j-i)$; the supplied
exact prime-exponent computation records each factor. Since a full
$V_n$ has rank $4^n$, even rational isomorphism to a whole $V_n$ is
impossible. The eightfold coefficient tensor retains the odd prime
factors. This does not obstruct the primitive embeddings or tensor
realizations constructed above, and says nothing about dimensions
of global invariant-form subspaces.

The per-embedding CM exponents of IF-10 are now
$(99,-205,-1085)$, with sums $(792,-1640,-8680)$ over all embeddings.
The same Tate-only obstruction and explicit CM-character untwist apply.
For the coordinate exponential family at these partitions, the positive
common-root construction of IF-09 fits inside $V_{39\cdot48}=V_{1872}$
per embedding, and inside $V_{14976}$ after tensoring all eight factors.

\paragraph{Exact scope.}
The input in (IF-11a) certifies a genuine algebraic theta-jet family.
It does not assert a nonzero composed terminal period, positive full
\texttt{sg}, or the full \texttt{SizeData}$_2$ hypotheses.
Indeed the separate source check SF-Q6-02 shows that the named witness
has top endpoint $b_T=-3$, incompatible with that full size condition.
No full step-10 divisibility bound is inferred from (IF-11e).
The movement here is that IF-08 and IF-10 now apply to an explicitly
traced, inhabited source theta family with the exact derivative
normalization, rather than to a hypothetical list of algebraic moments.
The algebraic geometry and $\ell$-adic translation are precisely those
of IF-10: specified coefficient correspondences and lattices, with the
CM character and all local denominators retained.
